\subsection{The construction}\label{sw:app:construction}

Except in the first subsection, whose statements hold for every channel and are used again in Appendix~\ref{sw:app:clock}, $\Phi$ is time-covariant. Throughout, $H$ and $\beta$ are fixed, entropies are in bits, and
\[
e_H(\rho)=\tfrac{\beta}{\ln2}\Tr H\bigl(\Phi(\rho)-\rho\bigr)=\Tr\rho\,D_H,\qquad D_H=\tfrac{\beta}{\ln2}\bigl(\Phi^*(H)-H\bigr).
\]
An \emph{exact dilation} $\Dil=(U,\tau)$ is a finite unitary $U$ on the system and a bath $B$ of dimension $R$ in state $\tau$, with $\Phi(\rho)=\Tr_BU(\rho\otimes\tau)U^*$. Let $\Gamma(\rho)=\Tr_SU(\rho\otimes\tau)U^*$ be its bath output and $f_{\Dil}(\rho)=S(\Gamma(\rho))-S(\tau)$.

\subsubsection{One dilation for all inputs}

For each full-rank $\rho$, the real-time entropy cost is the infimum over exact dilations,
\begin{equation}\label{sw:eq:pointwise}
s_\rho(\Phi)=\inf_{\Dil} f_{\Dil}(\rho),
\end{equation}
by the pointwise smoothing identity (Lemma~\ref{sw:lem:smoothing}). An exact dilation is an exact extension, which gives $\le$. Conversely, an $\varepsilon$-feasible extension lifts to a dilation at diamond error $\varepsilon/\lambda_{\min}(\rho)$ with the same objective, and a repair that changes the objective by a quantity tending to zero, independently of the extension's size, makes it exact.

\begin{lemma}[Minimax]\label{sw:lem:minimax}
$\Trt(\Phi)=\inf_{\Dil}\max_\rho\bigl[f_{\Dil}(\rho)+e_H(\rho)\bigr]$, the infimum over finite exact dilations.
\end{lemma}

\begin{proof}
Write $F_{\Dil}=f_{\Dil}+e_H$. Each $F_{\Dil}$ is continuous and concave on the state space, and a finite convex combination $\sum_jp_jF_{\Dil_j}$ is realized by the direct sum of the $\Dil_j$ controlled by a flag register in state $\sum_jp_j|j\rangle\langle j|$; the flag's entropy cancels in $f$. Let $a=\sup_\rho\inf_{\Dil}F_{\Dil}(\rho)$, which equals $\Trt(\Phi)$ by~\eqref{sw:eq:pointwise}, and fix $\delta>0$. Each $\rho$ has a dilation with $F_{\Dil}(\rho)<a+\delta$, which persists on an open neighbourhood. By compactness, finitely many $F_1,\dots,F_N$ satisfy $\min_jF_j(\rho)<a+\delta$ for every $\rho$. The set $\{x\in\mathbb R^N:x_j\le F_j(\rho)\ \forall j\ \text{for some}\ \rho\}$ is closed, convex and closed downwards, and excludes $(a+\delta,\dots,a+\delta)$. A separating hyperplane has a normal with nonnegative coordinates, which we normalize to weights $p_j$; then $\max_\rho\sum_jp_jF_j(\rho)<a+\delta$, and the flagged direct sum realizes $\sum_jp_jF_j$. Restricting to full-rank inputs changes nothing, since $|f_{\Dil}(\rho)-S(\rho)|\le\log d$ uniformly in $\Dil$, so concavity controls $F_{\Dil}$ near the boundary uniformly.
\end{proof}

Fix a dilation $\Dil$ with $g:=\max_\rho[f_{\Dil}(\rho)+e_H(\rho)]\le\Trt(\Phi)+\eta$, and write $\sigma=S(\tau)$ and $v=g+\sigma=\max_\rho[S(\Gamma(\rho))+e_H(\rho)]$.

\subsubsection{A joint first-order inequality}

Let $\rho_*$ maximize $S(\Gamma(\rho))+e_H(\rho)$ and $\zeta=\Gamma(\rho_*)$. The support $\mathcal S$ of $\zeta$ contains the support of every $\Gamma(\rho)$: mixing in weight $s$ outside $\mathcal S$ would raise the entropy by order $-s\log s$ and change everything else by $O(s)$. Put $L_B=-\log_2\zeta$ on $\mathcal S$ and $0$ on $\mathcal S^\perp$. The directional derivative of the concave maximand at $\rho_*$ gives $\Tr\Gamma(\rho)L_B+e_H(\rho)\le v$ for every $\rho$, that is,
\begin{equation}\label{sw:eq:firstorder}
\Gamma^*(L_B)+D_H\ \le\ v\,I.
\end{equation}
For $t\ge0$ let $U_t=e^{t\beta H/2}Ue^{-t\beta H/2}$ (a similarity, not a physical operation) and
\[
A(t)=\Tr_B\bigl[(I\otimes\tau^{1/2})\,U_t^*\,(I\otimes2^{tL_B})\,U_t\,(I\otimes\tau^{1/2})\bigr].
\]
Then $A(0)=I$, $A(t)\ge0$, and $A'(0)=\ln2\,\Gamma^*(L_B)+\beta(\Phi^*(H)-H)\le v\ln2\,I$ by~\eqref{sw:eq:firstorder}. Since all dimensions are fixed, Taylor's theorem gives $C_0,t_0>0$ with
\begin{equation}\label{sw:eq:qt}
A(t)\ \le\ q(t)\,I,\qquad q(t)=\exp\bigl(t\,v\ln2+C_0t^2\bigr),\qquad0\le t\le t_0.
\end{equation}

\subsubsection{An energy register}

Let the additive group generated by the energy gaps of $H$ have rank $r$, with rationally independent generators $\omega_1,\dots,\omega_r$, and write $E_e=E_*+a_e\cdot\omega$ with integer vectors $a_e$ for the eigenvalues of $H$, with spectral projectors $P_e$. On the lattice $\mathbb Z^r$ let $T_u|b\rangle=|b+u\rangle$ and $H_L|b\rangle=(E_0+b\cdot\omega)|b\rangle$. Define
\[
\mathcal L(X)=\sum_{e,f}(P_e\otimes I)\,X\,(P_f\otimes I)\otimes T_{a_f-a_e}.
\]
Orthogonality of the $P_e$ and $T_uT_v=T_{u+v}$ make $\mathcal L$ a unital $*$-homomorphism, so $V=\mathcal L(U)$ is unitary. The term $P_eXP_f$ changes the system energy by $E_e-E_f$ and moves the lattice by $a_f-a_e$, so $V$ conserves $H+H_L$ exactly; the bath cells have zero Hamiltonian. With $Q(t)=e^{-t\beta H_L/2}$ one checks $Q(t)VQ(t)^{-1}=\mathcal L(U_t)$, and hence, averaging over a fresh bath cell,
\begin{equation}\label{sw:eq:effect}
\Tr_B\Bigl[(I\otimes\tau^{1/2})\,V^*\bigl(I\otimes2^{tL_B}\otimes e^{-t\beta H_L}\bigr)V\,(I\otimes\tau^{1/2})\Bigr]=Q(t)\,\mathcal L(A(t))\,Q(t)\ \le\ q(t)\,I\otimes e^{-t\beta H_L}.
\end{equation}
This is an operator inequality on system and lattice together, so it holds for every joint input, with any reference attached. Only a finite box of the lattice is used: with $A_{\max}=\max_e\|a_e\|_\infty$, the box $\|b\|_\infty\le2(n+1)A_{\max}$ contains every coordinate reachable in $n$ rounds from $b=0$, the lifted unitary is specified exactly on reachable columns, and it is completed to a unitary separately inside each total-energy eigenspace. The box has dimension polynomial in $n$. The register starts in $|0\rangle$, a sharp energy eigenstate of energy $E_0=O(n)$ chosen to make all energies nonnegative, and it is counted as part of the battery.

\subsubsection{Why the register does not disturb the channel}

Fourier transformation diagonalizes the lattice translations, and $|0\rangle$ has the uniform distribution on the torus generated by the gaps. In the fibre of a phase $\theta$, the lifted unitary acts as $D_\theta UD_\theta^*$, where $D_\theta$ is the corresponding energy phase on the system. Tracing out the register therefore averages the whole adaptive experiment, with every round conjugated by the same $D_\theta$. Time covariance, with continuity, makes $\Phi$ commute with every $D_\theta$, so every fibre runs $n$ ideal uses of $\Phi$ and so does the average. This is the only place where time covariance is used.

\subsubsection{The weighted projector}

Run $n$ rounds, each with the lifted unitary on a fresh cell $B_i$ and the same register, and keep the used cells until the end. Let $W_j(t)=\prod_{i\le j}2^{tL_{B_i}}\otimes e^{-t\beta H_L}$. Before round $j$ the fresh cell is independent of everything else; the old factors are positive and act on other systems; and the observer's operations act on neither the cells nor the register. So~\eqref{sw:eq:effect} multiplies the expectation of $W_j(t)$ by at most $q(t)$ per round, whatever the correlations between system, register, old cells and references, and
\begin{equation}\label{sw:eq:supermart}
\Tr\Bigl[\rho_{\mathrm{fin}}\prod_{i=1}^n2^{tL_{B_i}}\,e^{-t\beta(H_L-E_0)}\Bigr]\ \le\ q(t)^n.
\end{equation}
Let $P$ be the spectral projector, inside $\mathcal S^{\otimes n}$, of $\sum_iL_{B_i}-\beta(H_L-E_0)/\ln2$ for values at most $n(v+\gamma)$. By Markov's inequality and~\eqref{sw:eq:supermart} with $t=\gamma\ln2/(2C_0)$, its rejected weight is at most $\exp[-n(\gamma\ln2)^2/(4C_0)]$ for small $\gamma$, for every observer. For each lattice label $b$, $P$ admits only bath eigenstrings with $\prod_iz_{j_i}\ge2^{-n(v+\gamma)}e^{-\beta(E_L(b)-E_0)}$, where $z_j$ are the eigenvalues of $\zeta$; there are at most $2^{n(v+\gamma)}e^{\beta(E_L(b)-E_0)}$ of them. Hence $P$ commutes with $H_L$ and
\begin{equation}\label{sw:eq:gibbsvolume}
Z_P:=\Tr P\,e^{-\beta H_L}\ \le\ D_L\,e^{-\beta E_0}\,2^{n(v+\gamma)},
\end{equation}
where $D_L$ is the dimension of the register. Counting the rank of $P$ instead of its Gibbs weight would lose exactly the cancellation between entropy and energy.

\subsubsection{Forming the baths}

Take $\gamma=\gamma_n=n^{-1/4}$. The controller $C$ consists of the $n$ bath cells, a residue register of $2^{r_n}$ levels with $r_n=\lceil3n\gamma_n\rceil+4$, and an $(n+1)$-level clock, all with zero Hamiltonian, and starts maximally mixed of dimension $m_C=(n+1)R^n2^{r_n}$. A unitary on $C$ maps a flat state of rank $K=2^{\lceil n(\sigma+2\gamma_n)\rceil}$ to a state whose cell marginal is within $2e^{-c_\tau n\gamma_n^2}+2^{-n\gamma_n}$ of $\tau^{\otimes n}$: assign $\lfloor Kp_x\rfloor$ flat labels to each typical eigenstring $x$ of $\tau^{\otimes n}$, keeping distinct residue labels for its whole fibre. The flat state is formed thermally by a standard construction. Let $W_f=\beta^{-1}\ln(m_C/K)$, and take a reservoir with levels $jW_f$, $j=0,\dots,L_f$, of integer degeneracies $m_C^jK^{L_f-j}$, so that every level has the same Gibbs weight, and a work qubit with energies $0,W_f$ starting excited. Swapping the equal-energy, equal-dimension subspaces $C\otimes|1\rangle\otimes R_j$ and $\mathrm{ran}(\Pi)\otimes|0\rangle\otimes R_{j+1}$ for $j<L_f$ is energy-conserving, and produces the flat state with the qubit in its ground state with probability $1-1/(L_f+1)$, without any measurement. Every later gate on $C$ conserves energy because $H_C=0$; the clock runs the schedule, so the whole implementation is one fixed unitary.

\subsubsection{Paying the work}

\begin{lemma}[Gibbs-weight converter]\label{sw:lem:converter}
Let $X$ have Hamiltonian $H_X$, let $P,Q$ be projectors commuting with $H_X$, and let $w<\beta^{-1}\ln(Z_Q/Z_P)$, with $Z_P=\Tr Pe^{-\beta H_X}$. For every $\delta>0$ there is a finite thermal reservoir and an energy-conserving unitary on $X$, the reservoir and a two-level work register with energy gap $w$ such that, with an input-independent probability at least $1-\delta$, every state supported on $P$, with any reference, is sent into $Q$ while the work register moves from its initial to its final level. On the remaining branch, $P$-supported inputs are left unchanged.
\end{lemma}

\begin{proof}
Write the energies of $X$ and the value $w$ in an integer basis of the group they generate: $\epsilon_0+a\cdot\omega$ and $w=w_{\mathrm v}\cdot\omega$. Let $p_a,q_c$ be the ranks of $P,Q$ in energy sector $a,c$. Give reservoir tile $j\in J=\{\|j\|_\infty\le L\}$ the energy $E_R^0+j\cdot\omega$ and the integer degeneracy $g_j=\lfloor Ne^{\beta j\cdot\omega}\rfloor$, with $N$ so large that $(1-\theta)Ne^{\beta j\cdot\omega}\le g_j$ on $J$, where $s=\beta^{-1}\ln(Z_Q/Z_P)-w>0$ and $\theta\le(1-e^{-\beta s})/2$. All tiles have nearly equal Gibbs weight, so the interior $G=\{\|j\|_\infty\le L-D\}$, $D=\max\|a-w_{\mathrm v}-c\|_\infty$, carries probability at least $1-\delta$ for $L$ large. For each total coordinate $u$, let $I_u$ be the source subspace $P\otimes R_G\otimes|i\rangle$ with $a+j=u$ and $J_u$ the target $Q\otimes R_J\otimes|f\rangle$ with $c+k+w_{\mathrm v}=u$. Then
\[
\frac{\dim J_u}{\dim I_u}\ \ge\ \frac{(1-\theta)Ne^{\beta(u\cdot\omega-w)}\sum_cq_ce^{-\beta c\cdot\omega}}{Ne^{\beta u\cdot\omega}\sum_ap_ae^{-\beta a\cdot\omega}}=(1-\theta)e^{\beta s}>1.
\]
Swap each $I_u$ with an isometric image inside $J_u$ and act as the identity elsewhere. Distinct shells are orthogonal, so this is an energy-conserving unitary, and on the reservoir block $G$ it moves the whole $P$-subspace, coherently, into $Q$ with the work register in $|f\rangle$.
\end{proof}

After the last output, apply the converter with $X=C\otimes L$, $H_X=H_L$, the cut $P_X=P\otimes I_{\mathrm{residue}}\otimes I_{\mathrm{clock}}$ and the target $Q_X=I_C\otimes|0\rangle\langle0|_L$, for which $Z_{Q_X}=m_Ce^{-\beta E_0}$. By~\eqref{sw:eq:gibbsvolume} one can extract
\[
w_n=\beta^{-1}\ln\frac{m_Ce^{-\beta E_0}}{(n+1)2^{r_n}D_L\,e^{-\beta E_0}\,2^{n(v+\gamma_n)}}-s_n
\]
with any fixed slack $s_n>0$, and the register returns to $|0\rangle$. Finally swap $C$ with a fresh maximally mixed cell of the same dimension, which returns the controller exactly and independently. On the good branch the net work is
\begin{align*}
W_f-w_n&=\kB T\ln2\,\bigl[n(v+\gamma_n)-\lceil n(\sigma+2\gamma_n)\rceil+r_n+\log_2(n+1)+\log_2D_L\bigr]+s_n\\
&=n\kB T\ln2\,\bigl[g+O(n^{-1/4})\bigr].
\end{align*}
No delivered output is touched, and no part of the device is reset during the $n$ rounds.

\subsubsection{The battery and the error}

The battery is the formation qubit, the converter's work register and the lattice register. It starts in one sharp energy eigenstate, its dimension is polynomial in $n$, and all its energies are within $O(n)$. On the good branch it ends in a fixed sharp eigenstate. The failure probability $q_n$ is at most the sum of the formation failure, the typicality error, the rejected weight of $P$ and the converter failure; with $L_f$ and the converter's reservoir chosen so that the formation and converter failures are at most $n^{-8}$, $q_n=O(n^{-8})$. Because the energy range is $O(n)$, the mean work differs from the good-branch value by $O(nq_n)$, the battery's energy variance and entropy tend to zero, and the battery ends nearly independent of everything else. The complete error of the experiment is at most the formation and typicality errors, since everything after the last output acts on internal systems. Choosing dilations with $g\le\Trt(\Phi)+\eta$ and letting $\eta\to0$ slowly with $n$ gives one family with vanishing error and work rate $\Trt(\Phi)$, which proves the upper half of Theorem~\ref{sw:thm:law}.

\subsection{The fixed-error converse}\label{sw:app:fixed-error}

Let $T=CR$ be the controller together with all reservoirs, with Hamiltonian $H_T$, initially in the Gibbs state $\gamma_T$ and independent of a battery $B$ in a pure eigenvector $|b_i\rangle$ of $H_B$ with energy $E_i$. For a projector $P_B$ commuting with $H_B$ write
\[
b_\beta(P_B)=\tfrac{\beta E_i}{\ln2}+\log_2\Tr\bigl[P_Be^{-\beta H_B}\bigr]=\log_2\frac{\Tr P_Be^{-\beta H_B}}{e^{-\beta E_i}},
\]
the logarithm of the Gibbs weight of $P_B$ relative to the initial battery state. For a target eigenvector of energy $E_f$, $b_\beta=\beta(E_i-E_f)/\ln2$ is the work in bits. We bound the error against one observer: it feeds every round half of a fresh purification $\psi_\rho$ of a fixed input $\rho$, keeps the references, applies a fixed test to a random fraction of the rounds, and at the end measures the total energy of its systems. Since this observer is not adaptive, the bound holds for the complete error.

\subsubsection{\texorpdfstring{An input that commutes with $H$, and an affine price}{An input that commutes with H, and an affine price}}

Let $\bar s_\rho=s_\rho(\Phi)$. By~\eqref{sw:eq:pointwise}, $\rho\mapsto\bar s_\rho$ is an infimum of concave functions and hence concave. Time covariance makes it invariant under $\rho\mapsto e^{-iHt}\rho e^{iHt}$, and the energy term is invariant and affine. Averaging over the closure of the time translations therefore does not decrease $\bar s_\rho+e_H(\rho)$, so for every $\eta>0$ there is a full-rank $\rho$ with $[\rho,H]=0$ and $\bar s_\rho+e_H(\rho)\ge\Trt(\Phi)-\eta$. Fix it.

On states $\nu$ of $Y=QO$ with $\nu_Q=\rho^T$, let $g(\nu)=\inf S(QO|F)$ over extensions with $\nu_{QF}=\rho^T\otimes\sigma_F$. Mixing extensions with a flag makes $g$ convex, and its lower semicontinuous hull at $\chi_\rho$ is $\bar s_\rho$. By separation there is an affine $\ell\le g$ with $\ell(\chi_\rho)\ge\bar s_\rho-\zeta$, for any $\zeta>0$. Give the reference the bookkeeping Hamiltonian $-H^T$ and set $H_Y=-H^T\otimes I+I\otimes H$; then $\psi_\rho$ has $H_Y$-energy zero and, by covariance, $[\chi_\rho,H_Y]=0$. Averaging $\ell$ over the group generated by $H_Y$ keeps it affine, below $g$ and unchanged at $\chi_\rho$, and makes it depend only on the pinching of $\nu$ into the eigenspaces of $H_Y$. The device never acts on $Q$; the reference Hamiltonian is only used in the proof.

\subsubsection{An energy-preserving test}

In each eigenspace of $H_Y$, of dimension $m_j$, take the rank-one informationally complete measurement obtained by normalizing the frame $\{e_i\}\cup\{e_i+e_k\}\cup\{e_i+\mathrm ie_k\}_{i<k}$; together these give effects $\Lambda_x=|v_x\rangle\langle v_x|$, $L=\sum_jm_j^2\le d^4$ of them, from whose statistics a fixed linear map recovers the pinching of $\nu$. Fix $0<\gamma<1/4$ and let the observer apply in each round the instrument with Kraus operators $\sqrt{1-\gamma}\,I_Y$ (no test) and $\sqrt\gamma\,|u_j\rangle\langle v_x|$, where $u_j$ is a fixed unit vector in the same eigenspace as $v_x$. Every Kraus operator commutes with $H_Y$, and the instrument has a dilation that preserves $H_Y$ exactly, with outcome registers of zero energy.

Let $A_i$ be the round-$i$ system $Y_i$ together with its outcome. For every state of the device before round $i$, purified by a register $Z$ that the device does not touch, the fresh input makes $Q_iZ$ a product, so $S(Q_iO_i|Z)\ge g(\nu_{Q_iO_i})\ge\ell(\nu_{Q_iO_i})$. On a tested outcome $Y_i$ is pure and independent of $Z$. Hence $S(A_i|Z)\ge(1-\gamma)\,\ell(\nu_{Q_iO_i})$, and since the outcome statistics $q$ determine the pinching of $\nu$, the affine function $f(q)=(1-\gamma)\,\ell(\text{pinching recovered from }q_x/\gamma)$ is a min-tradeoff function for every input. No dimension of the device enters $f$.

\subsubsection{Entropy accumulation}

Let $q_*$ be the ideal distribution of outcomes, fix a slack $\zeta_s>0$, put $h=f(q_*)-\zeta_s$, and let $G$ be the event that the empirical score is at least $h$ and at most $2\gamma n$ rounds are tested. The ideal experiment fails $G$ with probability exponentially small in $n$, and
\[
|G|\le\sum_{k\le2\gamma n}\binom nk L^k\le2^{n\theta_\gamma},\qquad\theta_\gamma=h(2\gamma)+2\gamma\log_2L.
\]
With trivial Markov conditions and side information the thermofield purifier $E$ of $\gamma_T$, which nothing touches, the entropy accumulation theorem~\cite{DFR20} gives, whenever $p_G=\Pr(G)>0$ and $0<\delta<1$,
\begin{equation}\label{sw:eq:eat}
H_{\min}^\delta(A^n|E)_{|G}\ >\ nh-C_f\sqrt{n\,[1-2\log_2(\delta p_G)]},
\end{equation}
with a constant $C_f$ depending only on $d$, $L$ and the gradient of $f$.

\subsubsection{The energy cut}

Let $\Delta E=\Tr H(\Phi(\rho)-\rho)$ and $u>0$. The instrument preserves $H_Y$, so in the ideal experiment the total energy $H_{A^n}=\sum_iH_{Y_i}$ is a sum of independent bounded terms with mean $n\Delta E$, and the projector $P_A$ onto $H_{A^n}\ge n(\Delta E-u)$ rejects with probability exponentially small in $n$. It is applied only at the end, in the analysis. Let $p$ be the probability that $G$, $P_A$ and the battery projector $P_B$ all succeed in the real experiment; these projectors commute. If $e_n$ bounds the ideal rejection, then
\begin{equation}\label{sw:eq:p}
p\ \ge\ 1-\varepsilon_n-e_n-q,
\end{equation}
where $q$ is the probability, in this experiment, that the battery ends outside $P_B$. Finitely many small $n$ are absorbed into the constant $c$.

\subsubsection{A Gibbs-weighted cap}

\begin{lemma}\label{sw:lem:cap}
Let $\omega_{AEK}$ be pure with $(H_A+H_K-H_E)\,\omega=E_i\,\omega$. Suppose $H_A\ge L_E$ on the support of $\omega_A$, $\omega_E\le e^{-\beta H_E}/(pZ_E)$, and $\omega_K$ is supported in a subspace commuting with $H_K$ of Gibbs weight $Z_K$. Then for $0<s<1$
\[
H^s_{\min}(A|E)_\omega\ \le\ \log_2\frac{e^{\beta(E_i-L_E)}Z_K}{pZ_E}+2\log_2\frac2{1-s}.
\]
\end{lemma}

\begin{proof}
Because the global state lies in one eigenspace of the total charge, $\omega_{AE}$ commutes with $H_A-H_E$, and in the block paired with $K$-energy $k$ its rank is at most the dimension $d_k$ of the $K$-support at that energy. Let $a_0=(1-s)/(2Z_K)$ and let $\Pi$ project onto the eigenvectors of $\omega_{AE}$ whose eigenvalue in block $k$ is at least $a_0e^{-\beta k}$. The rejected trace is at most $a_0\sum_kd_ke^{-\beta k}=(1-s)/2$, and $\omega_{AE}\ge a_0\,e^{\beta(H_A-H_E-E_i)}\,\Pi$, all three operators commuting. Using $e^{-\beta H_A}\le e^{-\beta L_E}$ on the support and tracing out $A$,
\[
\Tr_A\Pi\ \le\ \frac{e^{\beta(E_i-L_E)}}{a_0}\,e^{\beta H_E/2}\omega_Ee^{\beta H_E/2}\ \le\ \frac{e^{\beta(E_i-L_E)}}{a_0pZ_E}\,I_E.
\]
Any subnormalized $\sigma_{AE}$ within purified distance $s$ of $\omega_{AE}$ has $\Tr\Pi\sigma\ge(1-s)/2$. If $\sigma_{AE}\le2^{-h_0}I_A\otimes\zeta_E$ for a state $\zeta_E$, then
\[
\tfrac{1-s}2\ \le\ \Tr\Pi\sigma_{AE}\ \le\ 2^{-h_0}\Tr\bigl[(\Tr_A\Pi)\,\zeta_E\bigr]\ \le\ 2^{-h_0}\,\frac{e^{\beta(E_i-L_E)}}{a_0pZ_E}.
\]
Substituting $a_0$ and optimizing over $\sigma_{AE}$ and $\zeta_E$ gives the bound.
\end{proof}

Apply the lemma to the normalized state after the three projections, with $A=A^n$, $E$ the thermofield purifier of $\gamma_T$ (so $H_E=H_T^T$, $Z_E=Z_T$), and $K$ the device, the battery and the observer's outcome registers. Energy conservation of every gate and of the observer's dilation gives the charge condition. The two copies of the accepted outcome string span a space of dimension $|G|$ at energy zero, so $Z_K=Z_T\,\Tr[P_Be^{-\beta H_B}]\,|G|$, and the reservoir's partition function cancels:
\begin{equation}\label{sw:eq:cap}
H^{1/2}_{\min}(A^n|E)_{\mathrm{proj}}\ \le\ b_\beta(P_B)-\tfrac{\beta n(\Delta E-u)}{\ln2}+n\theta_\gamma+\log_2\tfrac1p+4.
\end{equation}
No bound on the size or spectrum of the reservoir is used.

\subsubsection{\texorpdfstring{From the event $G$ to all three projections}{From the event G to all three projections}}

\begin{lemma}\label{sw:lem:postproj}
Let $\psi_{AEK}$ be pure and normalized, let $F=P\otimes P_K$ with projectors $P$ on $A$ and $P_K$ on $K$, let $r=\|F\psi\|^2>0$ and $\delta\le r/64$. Then $H^{1/2}_{\min}(A|E)_{F\psi/\sqrt r}\ge H^{\delta}_{\min}(A|E)_\psi+\log_2r$.
\end{lemma}

\begin{proof}
Take a near-optimal smoothing $\sigma_{AE}$ within purified distance $\delta$ of $\psi_{AE}$ and, by Uhlmann's theorem, a pure lift $\phi$ with $\|\psi-\phi\|\le\sqrt2\delta$. The filtered marginal $\sigma'_0=\Tr_KF\phi F$ satisfies $\sigma'_0\le P\sigma_{AE}P$, so if $\sigma_{AE}\le2^{-h_0}I_A\otimes\zeta_E$ then also $\sigma'_0\le2^{-h_0}I_A\otimes\zeta_E$; and $\Tr\sigma'_0\le r+\delta$. The state $\sigma'=\sigma'_0/(r+\delta)$ is subnormalized and has conditional min-entropy at least $h_0+\log_2(r+\delta)\ge h_0+\log_2r$. Its fidelity with $F\psi/\sqrt r$ is at least $(\sqrt r-\sqrt2\delta)/\sqrt{r+\delta}$, so their squared purified distance is at most $4\delta/r\le1/16$, and $\sigma'$ is an admissible smoothing at level $1/2$.
\end{proof}

\subsubsection{Assembly}

Combine~\eqref{sw:eq:eat} with $\delta=p/64$ and $p_G\ge p$, Lemma~\ref{sw:lem:postproj} with $r\ge p$, and~\eqref{sw:eq:cap}:
\[
b_\beta(P_B)\ >\ n\Bigl[h+\tfrac{\beta(\Delta E-u)}{\ln2}-\theta_\gamma\Bigr]-C_f\sqrt{n\,[13+4\log_2(1/p)]}-2\log_2\tfrac1p-4.
\]
Choose $\eta,\zeta,\gamma,\zeta_s,u$ so small that the bracket is at least $\Trt(\Phi)-\delta'/2$, where $\delta'$ is the gap of Theorem~\ref{sw:thm:fixed-error}; all these choices are made before $n$. If $b_\beta(P_B)\le n(\Trt(\Phi)-\delta')$ and $\log_2(1/p)\le an$ with $a=\min\{\delta'/16,\delta'^2/(256C_f^2)\}$, the right side exceeds the left for $n$ large, a contradiction. So $p<2^{-an}$, and~\eqref{sw:eq:p} gives $1-\varepsilon_n\le q+e_n+2^{-an}$, which is Theorem~\ref{sw:thm:fixed-error}.

\subsubsection{Which battery contracts are enough}

\begin{remark}[The battery must be counted]\label{sw:rem:counted}
Sharpness of energy alone is not enough. At zero Hamiltonian, a battery made of $n$ pure qutrits with zero energy can be swapped into the controller to run the exact Werner--Holevo implementation at zero energy cost. Its entropy then rises by $n\log_23$. In Theorem~\ref{sw:thm:fixed-error} this is paid through $b_\beta(P_B)$, which counts the dimension of the accepted battery subspace; a projector of rank $R_B$ at energies at least $E_i-W$ has $b_\beta\le\beta W/\ln2+\log_2R_B$.
\end{remark}

\begin{remark}[Mean work]\label{sw:rem:wernermean}
For the Werner--Holevo channel at zero Hamiltonian, form the pure bath cells with a thermal transducer that succeeds with a fixed probability $p<1$ and run an arbitrary implementation on failure. The error is at most $1-p$, the battery entropy is $h(p)$, and the mean work per use tends to $p\log_23<\Trt$. The theorem is not contradicted: the battery misses its target with probability $1-p$, and the work on the good branch is at the full rate.
\end{remark}

\subsection{What is free in real time}\label{sw:app:zero-set}

For a dilation $\Dil=(U,\tau)$ of a channel $\Psi$, let $\Omega=U(\gamma\otimes\tau)U^*$, $\omega=\Tr_S\Omega$, and define the \emph{Gibbs defect}
\[
G_{\Dil}=S(\omega)-S(\tau)+\tfrac{\beta}{\ln2}\Tr H\bigl(\Psi(\gamma)-\gamma\bigr)=f_{\Dil}(\gamma)+e_H^\Psi(\gamma).
\]
Since $S(\Omega)=S(\gamma)+S(\tau)$ and $\log_2\gamma=-\beta H/\ln2-\log_2Z$, expanding the relative entropy gives
\begin{equation}\label{sw:eq:defect}
G_{\Dil}=D\bigl(\Omega\,\big\|\,\gamma\otimes\omega\bigr)\ \ge\ 0,
\end{equation}
with no assumption on $\tau$. For a thermal operation, $U$ commutes with $\gamma\otimes\tau$, so $G_{\Dil}=0$.

\subsubsection{From zero cost to dilations with vanishing defect}

Since $s_\rho(\Phi)\ge S(\rho)-S(\Phi(\rho))$, $\Trt(\Phi)\ge D(\Phi(\gamma)\|\gamma)$; so $\Trt(\Phi)=0$ forces $\Phi(\gamma)=\gamma$ and $s_\gamma(\Phi)=0$. Take extensions at input $\gamma$ with errors $\varepsilon_j\to0$ and objectives $u_j\to0$. Purifying such an extension and using uniqueness of purification, as in Appendix~\ref{sw:app:construction}, turns it into a finite dilation $\Dil_j$ with $f_{\Dil_j}(\gamma)=u_j$, whose channel $\Psi_j$ is within $\varepsilon_j/\gamma_{\min}$ of $\Phi$ in half diamond norm. Since $\Phi(\gamma)=\gamma$, the energy term of $\Psi_j$ at $\gamma$ is at most $\beta\,\mathrm{span}(H)\varepsilon_j/(\gamma_{\min}\ln2)$, so $G_{\Dil_j}\to0$. Only the input $\gamma$ is used. This proves the single-input criterion stated after Theorem~\ref{sw:thm:free}.

\subsubsection{A small defect forces time covariance}

For states $\rho,\sigma$ with eigenpairs $(r_i,|i\rangle)$ and $(s_j,|j\rangle)$, the distributions $p_{ij}=r_i|\langle i|j\rangle|^2$ and $q_{ij}=s_j|\langle i|j\rangle|^2$ satisfy $D(p\|q)=D(\rho\|\sigma)$, and Jensen's inequality gives $\Tr\sqrt\rho\sqrt\sigma\ge e^{-D_{\mathrm{nat}}(\rho\|\sigma)/2}$, hence
\begin{equation}\label{sw:eq:sqrtdist}
\|\sqrt\rho-\sqrt\sigma\|_2^2\le D_{\mathrm{nat}}(\rho\|\sigma).
\end{equation}
Add a fresh copy $A$ of the system in state $\gamma$ to the bath, let $D=AB$, $\sigma_D=\gamma_A\otimes\tau_B$ and $W=U_{AB}^*U_{SB}$. The second factor acts on the bath only, so $W$ with bath state $\sigma_D$ implements exactly the same channel $\Psi$. Applying~\eqref{sw:eq:defect} and~\eqref{sw:eq:sqrtdist} once on $SB$ and once on $AB$,
\begin{equation}\label{sw:eq:commutator}
\bigl\|[W,\sqrt{\gamma\otimes\sigma_D}]\bigr\|_2\ \le\ 2\sqrt{\ln2\,G_{\Dil}}.
\end{equation}
Replace $\sigma_D$ by $\sigma_t=(1-t)\sigma_D+tI/\dim D$; this changes the channel by at most $t$ and the commutator by at most $2\sqrt{2t}$. For $\eta=\gamma\otimes\sigma_t$ and real $z$, the elementary bound $|x^{iz}-1|\le C_z|\sqrt x-1|$, with $C_z=\max\{2e^{1/2}|z|,\,2/(1-e^{-1/2})\}$, applied to matrix elements gives $\|[W,\eta^{iz}]\|_{2,\eta}\le C_z\|[W,\sqrt\eta]\|_2$. The unitary $\eta^{-iz}W\eta^{iz}$ implements the channel conjugated by $\gamma^{iz}$, since the bath factor disappears under the bath trace, and comparing Stinespring vectors weighted by $\eta$ gives a diamond distance at most $C_z\|[W,\sqrt\eta]\|_2/\sqrt{\gamma_{\min}}$. If $G_{\Dil_j}\to0$ and $\Psi_j\to\Phi$, then $\Phi$ commutes with conjugation by $\gamma^{iz}$ for every $z$. Since $\gamma^{iz}$ is $e^{-i\beta zH}$ up to a phase, $\Phi$ is time-covariant.

\subsubsection{Twirling preserves the defect}

Let $H$ have $r$ rationally independent energy gaps, with $E_e=E_*+a_e\cdot\omega$ for integer vectors $a_e$, and $D_\theta=\sum_ee^{\mathrm ia_e\cdot\theta}P_e$. A channel is time-covariant exactly when it commutes with every $D_\theta$, because $E_e-E_f=E_g-E_h$ if and only if $a_e-a_f=a_g-a_h$. Averaging over $\theta\in(2\pi/Q)\{0,\dots,Q-1\}^r$, with $Q$ larger than the absolute value of every coordinate of $a_e-a_f-a_g+a_h$, implements the full twirl exactly, and is realized by a uniform flag in the bath controlling the conjugations. At input $\gamma$ every conjugated dilation has the same bath output, the flag's entropy cancels, and the energy of the system is unchanged, so the defect is exactly preserved. We may therefore assume the dilations $\Dil_j$ have time-covariant channels.

\subsubsection{Replacement}

\begin{theorem}[Replacement]\label{sw:thm:replacement}
There are constants $c_{H,\beta},G_0>0$, depending only on $H$ and $\beta$, such that every dilation $\Dil$ of a time-covariant channel $\Psi$ with $0<G_{\Dil}\le G_0$ admits a finite thermal operation $\Theta$ at the same $H$ and $\beta$ with
\[
\tfrac12\|\Psi-\Theta\|_\diamond\ \le\ c_{H,\beta}\,G_{\Dil}^{1/(6\max\{1,r\})}.
\]
\end{theorem}

The constant does not depend on the dimension of the original bath or on its smallest eigenvalue. Together with the two previous subsections, Theorem~\ref{sw:thm:replacement} proves the hard direction of Theorem~\ref{sw:thm:free}; the easy direction is in Section~\ref{sw:sec:free}. We give the proof in four steps.

\emph{Commensurate energies.} Suppose first that every gap of a Hamiltonian $H_0$ is an integer multiple of $\Delta_0$, with Gibbs state $\gamma_0$ and $m_0=\lambda_{\min}(\gamma_0)$. Choose a mesh $h_E$ dividing $\Delta_0$, put $h=\beta h_E$, $\eta=\gamma_0\otimes\sigma_t$ and $q=\|[W,\sqrt\eta]\|_2^2$. In an eigenbasis of $\eta$ with eigenvalues $p_a$, $q=\sum_{ab}|W_{ab}|^2(\sqrt{p_a}-\sqrt{p_b})^2$. Give each bath eigenvector the logarithmic energy $\xi_b=-\ln\sigma_{t,b}$, pick an offset $s$ uniformly in $[0,h)$, and round $\xi_b$ down to the grid $h\mathbb Z-s$. The rounded Gibbs state $\sigma'$ satisfies $e^{-h}\sigma_t\le\sigma'\le e^h\sigma_t$. Because the gaps of $H_0$ are multiples of $h_E$, the rounded total-energy blocks are exactly the bins of $-\ln p_a$. Two logarithms fall in different bins with probability $\min(1,|\ln(p_a/p_b)|/h)\le|\sqrt{p_a/p_b}-1|/(1-e^{-h/2})$, so by Cauchy--Schwarz the expected $\eta$-weighted mass $M_s$ of $W$ outside the diagonal blocks is at most $\sqrt q/(1-e^{-h/2})$. Fix an offset achieving this. In each block, complete $P_lWP_l$ to a unitary by its polar decomposition; the direct sum $V$ commutes exactly with $H_0$ plus the rounded bath Hamiltonian, and $\|W-V\|^2_{2,\eta'}\le2e^hM_s$ with $\eta'=\gamma_0\otimes\sigma'$. Since every input is dominated by $\gamma_0/m_0$, the thermal channel $\Theta_0$ of $(V,\sigma')$ satisfies
\[
\tfrac12\|\Psi-\Theta_0\|_\diamond\ \le\ t+\tfrac{e^h-1}2+\Bigl[\frac{2e^h\sqrt q}{m_0(1-e^{-h/2})}\Bigr]^{1/2}.
\]
No dimension or smallest eigenvalue of the bath appears.

\emph{Irrational energies.} Write $\omega=\Delta(1,v_2,\dots,v_r)$. By Dirichlet's theorem, for each integer $L$ there are $1\le q_L\le L^{r-1}$ and integers $p_k$ with $|q_Lv_k-p_k|\le1/L$. Let $\omega_0=\Delta(1,p_2/q_L,\dots,p_r/q_L)$ and $H_0=\sum_e(E_*+a_e\cdot\omega_0)P_e$, with mesh $h_E=\Delta/(mq_L)$, $m=\lceil L^{1/3}\rceil$. Then $\|H_0-H\|=O(1/(Lq_L))$, the original integer relations between gaps are preserved, and $\|\sqrt{\gamma_0}-\sqrt\gamma\|_2=O(\|H_0-H\|)$. The commutator bound~\eqref{sw:eq:commutator} enters linearly, so $\sqrt q\le2\sqrt{\ln2\,G_{\Dil}}+2\sqrt{2t}+O(1/(Lq_L))$, and with $t=G_{\Dil}$ the previous step gives $\tfrac12\|\Psi-\Theta_0\|_\diamond\le c_{H,\beta}[L^{-1/3}+G_{\Dil}^{1/4}L^{r/2-1/3}]+G_{\Dil}$.

\emph{Removing extra coincidences.} The approximation $H_0$ may have gap coincidences that $H$ lacks. Twirling $\Theta_0$ over the original grading, as above, gives a thermal operation $\Xi_0$ at $H_0$, since each $D_\theta$ commutes with $H_0$; and since $\Psi$ is covariant for $H$, twirling does not increase the distance to $\Psi$.

\emph{Conserving the original Hamiltonian exactly.} Add a lattice register $|\ell\rangle$, $\ell\in[-N,N]^r$, with Hamiltonian $E_L^0+\ell\cdot(\omega-\omega_0)$ in its Gibbs state. For $k$ in the inner box $[-N+A,N-A]^r$, $A=\max_e\|a_e\|_\infty$, the isometries $J_k(|e\rangle\otimes|b\rangle)=|e,b,k-a_e\rangle$ have orthogonal images, each reducing $H+H_E+H_L$, with $J_k^*(H+H_E+H_L)J_k=H_0+H_E+\text{const}$. So $V=\bigoplus_kJ_kV_0J_k^*\oplus I$ conserves the original total Hamiltonian exactly, where $V_0$ is the thermal unitary of $\Xi_0$. For initial labels in the box $[-N+2A,N-2A]^r$, tracing out the lattice reproduces exactly the original-grading twirl of $\Xi_0$, which is $\Xi_0$. The remaining labels have Gibbs probability at most $4rA/(2N+1)+2\beta N\|\omega-\omega_0\|_1$. With $N\sim L^{1/3}$ and $L=\lceil G_{\Dil}^{-1/(2r)}\rceil$, all terms are $O(G_{\Dil}^{1/(6r)})$. For $r\le1$, $H_0=H$ and the lattice is not needed, which gives the exponent $1/6$. The new bath has dimension $O(d\,R\,G_{\Dil}^{-1/6})$ if the original has dimension $R$.

\subsubsection{The qutrit channel}

\begin{proof}[Proof of Proposition~\ref{sw:prop:qutrit}]
Write $a=\sqrt{1-\lambda}$ and $(g_0,g_1,g_2)$ for the Gibbs populations, so $g_2=\lambda g_0$. The channel is covariant because each Kraus operator moves a definite energy, and Gibbs-preserving because its population map fixes $(g_0,g_1,g_2)$. Its minimal Stinespring isometry is $|0\rangle\mapsto a|0\rangle e_0+\sqrt\lambda|2\rangle e_+$, $|1\rangle\mapsto|1\rangle e_0$, $|2\rangle\mapsto|0\rangle e_-$. Any exact dilation, with its bath state purified by a register $R$ that the dilation does not touch, has this form with orthonormal vectors $v,z,w$ in place of $e_0,e_+,e_-$. Because $R$ is untouched, the map $X\mapsto\Tr_{SB}(\cdot)$ onto $R$ sends every $X$ to $\Tr(X)\,\tau_R$. Applied to $X=|0\rangle\langle2|$ it gives $a\Tr_B|v\rangle\langle w|=0$, so the reduced bath states $\nu_v,\nu_w$ of $v,w$ have fidelity $\|\Tr_B|v\rangle\langle w|\|_1=0$: they are orthogonal. At input $\gamma$, the bath output is $\omega=(1-2g_2)\nu_v+g_2\nu_z+g_2\nu_w$, and the output is $1$ exactly when the bath is in $\nu_v$, with probability $g_1$. Dephasing the system in~\eqref{sw:eq:defect} and keeping only the output-$1$ term,
\[
G_{\Dil}\ \ge\ g_1\,D(\nu_v\|\omega)\ \ge\ g_1\log_2\frac1{\Tr P\omega}\ \ge\ g_1\log_2\frac1{1-g_2},
\]
where $P$ is the support projector of $\nu_v$, orthogonal to $\nu_w$.

For an extension at error $\varepsilon$, lift it to a dilation as above, with channel within $\varepsilon'=\varepsilon/g_{\min}$ of $\Phi$. By the continuity of Stinespring dilations~\cite{KSW08}, after enlarging the bath it is within $t=\sqrt{2\varepsilon'}$ in operator norm of an exact dilation of $\Phi$ on the same untouched $R$. The argument above then holds approximately: $\|\Tr_B|v\rangle\langle w|\|_1\le2t/a$, a binary test separates $\nu_v$ from $\nu_w$ up to $2t/a$, and the mutual information between the event ``output $1$'' and that test is at least $g_1d_2(2t/a\,\|\,g_2(1-2t/a))-t\log_23-3h(t)$, where $d_2$ is the binary relative entropy. This tends to $g_1\log_2\frac1{1-g_2}$ as $\varepsilon\to0$, uniformly in the bath, and so does the energy correction $\beta\,\mathrm{span}(H)\varepsilon'/\ln2$ between the defect and the objective of the extension. Hence $s_\gamma(\Phi)\ge g_1\log_2\frac1{1-g_2}$, and $\Trt(\Phi)$ is at least this, since the energy term vanishes at $\gamma$. The same estimate with Pinsker's inequality gives $G_{\Dil}\ge\frac2{\ln2}[u-c\sqrt{2\delta}]_+^2$ for every dilation at diamond distance $\delta$ from $\Phi$, with $u=g_1g_2$ and $c=2+6u/a$. Since thermal operations have $G_{\Dil}=0$, every thermal operation is at least $u^2/(2c^2)$ from $\Phi$.
\end{proof}

\subsubsection{A single input}

Let $\Theta$ be a thermal operation whose bath has Hamiltonian $H_B=\sum_EE\,\Pi_E$, eigenspaces $\mathcal H_E$ of dimension $d_E$, and Gibbs weight $p_E=e^{-\beta E}/Z_B$ per state, so that $\sum_Ep_Ed_E=1$ and $p_E=\lambda p_{E-E_2}$. Write $T_{ji}=\langle j|\Theta(|i\rangle\langle i|)|j\rangle$. Since $U$ conserves energy, it maps $|i\rangle\otimes\mathcal H_E$ into states of total energy $E_i+E$. We use three of its blocks, all contractions: $A_E$ and $B_E$, the blocks of $U$ from $|0\rangle\otimes\mathcal H_E$ to $|0\rangle\otimes\mathcal H_E$ and from $|1\rangle\otimes\mathcal H_E$ to $|1\rangle\otimes\mathcal H_E$, and $C_E$, the block from $|0\rangle\otimes\mathcal H_E$ to $|2\rangle\otimes\mathcal H_{E-E_2}$. Expanding the partial trace and grouping bath states by energy,
\begin{equation}\label{sw:eq:modes}
\begin{gathered}
c:=\langle0|\Theta(|0\rangle\langle1|)|1\rangle=\sum_Ep_E\Tr(B_E^*A_E),\\
T_{00}=\sum_Ep_E\|A_E\|_2^2,\quad T_{11}=\sum_Ep_E\|B_E\|_2^2,\quad T_{20}=\sum_Ep_E\|C_E\|_2^2.
\end{gathered}
\end{equation}
By time covariance, $\langle0|\Theta(|i\rangle\langle k|)|1\rangle=0$ unless $E_k-E_i=E_1$, so for every input $\rho$ supported on levels $0$ and $1$ the $0$--$1$ coherence of $\Theta(\rho)$ is $c\,\rho_{01}$.

\begin{lemma}[Coherence and transfer]\label{sw:lem:coherence}
For every thermal operation, $T_{20}\le\lambda$ and
\[
|c|\ \le\ 1-\lambda+\sqrt{\lambda(\lambda-T_{20})}.
\]
The same holds for every limit of thermal operations.
\end{lemma}

\begin{proof}
By von Neumann's trace inequality and $\|B_E\|\le1$, $|\Tr(B_E^*A_E)|\le\|A_E\|_1$. The columns of $U$ on $|0\rangle\otimes\mathcal H_E$ are orthonormal and have components only at levels $0$, $1$ and $2$, so $A_E^*A_E\le I-C_E^*C_E$. The square root is operator monotone, hence
\[
\|A_E\|_1\ \le\ \Tr\sqrt{I-C_E^*C_E}\ =\ d_E-\sum_if(x_{E,i}),\qquad f(x)=1-\sqrt{1-x},
\]
where the $x_{E,i}$ are the squared nonzero singular values of $C_E$; there are at most $\operatorname{rank}C_E\le d_{E-E_2}$ of them. Give each pair $(E,i)$ the weight $p_E$. The total weight is
\[
W\ \le\ \sum_Ep_Ed_{E-E_2}\ =\ \lambda\sum_Ep_{E-E_2}d_{E-E_2}\ \le\ \lambda,
\]
and the weighted sum of the $x_{E,i}$ is $T_{20}$. Since each $x_{E,i}\le1$, $T_{20}\le W\le\lambda$. The function $f$ is convex with $f(0)=0$, so Jensen's inequality gives $\sum_{E,i}p_Ef(x_{E,i})\ge Wf(T_{20}/W)$, and $Wf(T_{20}/W)$ does not increase with $W$ because $f(z)\le zf'(z)$. Hence $\sum_{E,i}p_Ef(x_{E,i})\ge\lambda f(T_{20}/\lambda)=\lambda-\sqrt{\lambda(\lambda-T_{20})}$. Summing the bounds on $|\Tr(B_E^*A_E)|$ over $E$ with $\sum_Ep_Ed_E=1$ gives the claim. Both $c$ and $T_{20}$ are continuous functions of the channel and the bound does not depend on the bath, so it passes to limits.
\end{proof}

The lemma is tight. Take as bath a truncated oscillator with levels $0,E_2,\dots,KE_2$, one state each. In each total-energy shell, rotate $|0\rangle\otimes|kE_2\rangle$ into $|2\rangle\otimes|(k-1)E_2\rangle$ by a fixed angle with cosine $a$, and leave level $1$ alone. Then $c=1-W+Wa$ and $T_{20}=W(1-a^2)$, where $W=\lambda(1-\lambda^K)/(1-\lambda^{K+1})\to\lambda$, so the bound is attained in the limit $K\to\infty$ for every $T_{20}\in[0,\lambda]$. At $T_{20}=\lambda$ it bounds the coherence of $\Theta(\psi_0)$ by $(1-\lambda)/2$. A thermal operation whose output on $\psi_0$ has exactly the populations of $\Phi(\psi_0)$ has $T_{20}=\lambda$: the Gibbs-preservation equations force $T_{01}=0$ and $T_{02}=1$, and then $T_{20}=\lambda$. So the lemma recovers the exact value of Ding, Ding and Hu~\cite{DingFengHu21} without their assumption on the bath degeneracies.

\begin{proof}[Proof of the single-input bound in Proposition~\ref{sw:prop:qutrit}]
Let $\Theta$ be a limit of thermal operations, $X=\Theta(\psi_0)-\Phi(\psi_0)$ and $\varepsilon=\frac12\|X\|_1$. Write $s=\sqrt{1-\lambda}$. Then $\Phi(\psi_0)$ has populations $\bigl((1-\lambda)/2,\,1/2,\,\lambda/2\bigr)$ and $0$--$1$ coherence $s/2$, and by covariance $\Theta(\psi_0)$ has populations $p_j=(T_{j0}+T_{j1})/2$ and $0$--$1$ coherence $c/2$. Put $y=s-|c|$ and $x=\lambda-T_{20}\ge0$.

\emph{Step 1.} Since $X$ is traceless, $\|X\|\le\varepsilon$, so $T_{00}\le2p_0\le1-\lambda+2\varepsilon$. Pinching $X$ onto the blocks $\{0,1\}$ and $\{2\}$ does not increase its trace norm, and a $2\times2$ Hermitian matrix has trace norm at least twice the modulus of its off-diagonal entry. Hence
\[
2\varepsilon\ \ge\ |c-s|+\bigl|\tfrac\lambda2-p_2\bigr|\ \ge\ |y|+|e|,\qquad e=\tfrac12(\lambda-T_{20}-T_{21})=\tfrac12(x-T_{21}).
\]

\emph{Step 2.} By the Cauchy--Schwarz inequality in~\eqref{sw:eq:modes}, $|c|^2\le T_{00}T_{11}$, and $T_{11}\le1-T_{21}$. So $(s-y)^2\le(s^2+2\varepsilon)(1-T_{21})$, which gives $T_{21}\le(2\varepsilon+2sy)/(1-\lambda)$ if $y\ge0$ and $T_{21}\le2\varepsilon/(1-\lambda)$ if $y<0$.

\emph{Step 3.} From Step 1, $x=2e+T_{21}\le2(2\varepsilon-|y|)+T_{21}$. If $y\ge0$, then using $y\le2\varepsilon$ and $s/(1-\lambda)=1/s$,
\[
x\ \le\ \varepsilon\Bigl(4+\frac2{1-\lambda}\Bigr)+2y\Bigl(\frac1s-1\Bigr)\ \le\ \varepsilon\Bigl(\frac2{1-\lambda}+\frac4s\Bigr),
\]
and the same bound holds if $y<0$. So $\lambda x\le k\varepsilon$ with $k=2\lambda(1+2s)/(1-\lambda)$.

\emph{Step 4.} Lemma~\ref{sw:lem:coherence} gives $\sqrt{\lambda x}\ge|c|-(1-\lambda)=\Delta_\lambda-y\ge\Delta_\lambda-2\varepsilon$. If $2\varepsilon<\Delta_\lambda$, squaring gives $(\Delta_\lambda-2\varepsilon)^2\le k\varepsilon$, so $\Delta_\lambda^2\le(4\Delta_\lambda+k)\varepsilon$ and $\varepsilon\ge\Delta_\lambda^2/(4\Delta_\lambda+k)=\delta(\lambda)$. If $2\varepsilon\ge\Delta_\lambda$, then $\varepsilon\ge\Delta_\lambda/2\ge\delta(\lambda)$. The diamond-norm statement follows because $\frac12\|\Theta-\Phi\|_\diamond\ge\varepsilon$.
\end{proof}

The constant behaves as $\delta(\lambda)\approx\lambda/32$ as $\lambda\to0$ and $\delta(\lambda)\approx(1-\lambda)^2/2$ as $\lambda\to1$. The bound $u^2/(2c^2)$ is larger only at very high temperature, $\lambda\gtrsim0.97$. For equally spaced levels, $E_2=2E_1$, with $\beta E_1=\ln2$, so that $\lambda=1/4$, the bounds read $\Trt(\Phi)\ge\frac27\log_2\frac76\approx0.064$ and $\delta(1/4)\approx5.9\times10^{-3}$.

\subsection{A clock that returns approximately}\label{sw:app:clock}

The lower half of Theorem~\ref{sw:thm:clock} is~\eqref{sw:eq:converse}. Its proof uses a product initial state, energy conservation, real-time release, the complete error and the return of the controller's free energy up to $o(n)$, all of which hold when the controller returns approximately; it uses neither time covariance nor exact return. This appendix gives the construction and the bound on the baths' range.

\subsubsection{The rare-excitation primitive}

Let $U$ be a unitary on the system, $Q=U^*HU$, $D=Q-H$, $w=\lambda_{\max}(D)$, $B=wI-D\ge0$ and $b=\|B\|$. For $M>b$ let $G$ be a qubit with $H_G=M|1\rangle\langle1|$ in its Gibbs state $\gamma_G=(1-z)|0\rangle\langle0|+z|1\rangle\langle1|$, $z=(1+e^{\beta M})^{-1}$. Put $F=B/M$, $A=\sqrt{I-F}$ and $C=\sqrt F$; they commute and $A^2+C^2=I$. In the basis $|0\rangle,|1\rangle$ of $G$ define the unitary
\[
W_M=(U\otimes I)\begin{pmatrix}A&-C\\C&A\end{pmatrix}.
\]
For either input of $G$ the payload channel is $\Phi_M(\rho)=U(A\rho A+C\rho C)U^*$, and since $\|A-I\|\le b/M$ and $\|C\|^2=b/M$, $\frac12\|\Phi_M-U(\cdot)U^*\|_\diamond\le3b/(2M)$. The excitation probability of $G$ after the contact is $(1-z)F+z(I-F)$, so
\[
\Tr_G\gamma_G\bigl[W_M^*(H+H_G)W_M-(H+H_G)\bigr]=wI+R_M,
\]
with $R_M=AQA+CQC-Q-2zB$ and $\|R_M\|\le3\|H\|b/M+2zb$.
The input-dependent part $-D$ of the energy change has gone into the bath as the heat effect $(1-2z)B$, and what remains is the constant $w$ up to $R_M$. The heat leaves a record, but a small one: purifying $\gamma_G$, the contact maps $\psi$ to $UA\psi\otimes g_z+UC\psi\otimes h_z$ with orthonormal $g_z,h_z$, at squared distance at most $2b/M$ from $U\psi\otimes g_z$.

\subsubsection{Lift and pump}

Write the energies of the system and $G$ as integer combinations $a_x\cdot\omega$ of a basis $\omega$ of the group they generate together with $w$, and let a reference $P$ on $\mathbb Z^r$ have Hamiltonian $\omega\cdot J$ and translations $T_a$. The map $L(Y)=\sum_{x,y}\Pi_xY\Pi_y\otimes T_{a_y-a_x}$ is a unital $*$-homomorphism, so $L(W_M)$ is unitary, and it conserves the energy of system, $G$ and reference exactly. At each use, also translate the reference's $w$-coordinate by one while a battery label moves from $k$ to $k+1$, with battery energies $(n-k)w$. Because the initial state of $G$ is diagonal, bra and ket share the label of $G$, and
\[
\Tr_G\bigl[(I\otimes\gamma_G^{1/2})L(Y)(I\otimes\gamma_G^{1/2})\bigr]=L_S\bigl(\Tr_G\gamma_GY\bigr)
\]
for every $Y$, with $L_S$ the corresponding map on the system alone. Applied to the energy, this shows that the reference's energy changes by $-L_S(R_M)$ per use, an operator of norm at most $\|R_M\|$, whatever the joint state of input and reference.

\subsubsection{The Hadamard gate}

Let $H=E|1\rangle\langle1|$ and $U=(X+Z)/\sqrt2$, so $w=E/\sqrt2$ and $b=\sqrt2E$. Take $m=(n+1)^4+2$ and $M=mE$. The reference has two coordinates, with energies $E$ and $w$, in a finite box. It starts in a product of sine packets, the reference states Chiribella, Yang and Renner use for the same purpose~\cite{ChiribellaYangRenner21}, of widths $L_1=(n+1)^8$ and $L_2=(n+1)^4$, with buffers $(n+1)(m+1)+2$ and $n+2$, so that every history, including every rare heat branch, stays inside the box. On the box each lifted contact is completed to a square energy-conserving unitary within each energy shell. A counter with $n+1$ states and zero Hamiltonian schedules the contacts and is returned by a known permutation. The $n$ qubits $G$ are fresh, independent Gibbs stocks and are never reused.

In the Fourier representation of the reference, one common phase governs all contacts. The purified history $\Psi_\theta$ satisfies $\|\partial_1\Psi_\theta\|\le2n(m+1)$, $\|\partial_1^2\Psi_\theta\|\le4n^2(m+1)^2$ and $\|\partial_2\Psi_\theta\|\le n$, uniformly over observers, and the Fourier density of a sine packet of width $L$ is even with second moment at most $\pi^4/(4(L+1)^2)$. Hence the actual global state is within
\[
\delta_n=\pi^2\Bigl[\frac{n(m+1)}{L_1+1}+\frac n{L_2+1}\Bigr]=O(n^{-3})
\]
of the initial packet tensored with the history at zero phase, and the controller's joint state with all other systems is within $2\delta_n$ of the product. Evenness cancels the first-order phase effect for the observer, and at zero phase each use applies $\Phi_M$, so the error is at most $\frac32nb/M+\pi^4n^2(m+1)^2/(L_1+1)^2=O(n^{-3})$. The controller's energy changes by at most $n\|R_M\|=o(1)$ in mean and its entropy by $O(n^{-6}\log n)$. The battery starts at label $0$ and ends at label $n$ on every path, so the work is exactly $nE/\sqrt2$. The retired qubits absorb a total heat between $0$ and $2nw$, rarely and in large quanta.

\subsubsection{Every channel}

Fix an exact dilation $(U,\tau)$ of $\Phi$ with a bath of dimension $R$ and cost $g=\max_\rho[f(\rho)+e_H(\rho)]$, let $\zeta$ be the bath output at a maximizing input, and $v=g+S(\tau)$. The support of $\zeta$ contains that of every bath output, and the tangent at the maximizer gives the operator inequality
\[
\Gamma^*(-\log_2\zeta)+D_H\ \le\ vI.
\]
For $\delta\in(0,1)$ put $\zeta_\delta=(1-\delta)\zeta+\delta I/R$ and give the bath the Hamiltonian $H_B=-\beta^{-1}\ln\zeta_\delta$, whose partition function is one. With $V=(\ln2/\beta)v-\beta^{-1}\ln(1-\delta)$, the operator $VI-[\Phi^*(H)-H+\Gamma^*(H_B)]$ is positive. This turns the bath's entropy into energy: the primitive above, applied to the contact $U$ on the system and a fresh bath copy, absorbs the input-dependent energy with one positive operator, and a constant $w_c=V-\Tr\tau H_B$ is pumped per use.

The bath copies must start in $\tau^{\otimes n}$, which need not commute with $H_B$. Forming them, and returning everything used to form them, is the finite-source module of Appendix~\ref{app:module}, which uses nothing from this part, so the dependency is one-way. Its hypotheses hold here with a single slot: the retained body is one bath copy with Hamiltonian $H_B$, so that its faithful state is $q=\zeta_\delta$; there is one source per use; the gate is the contact $W_M$ of the primitive, with $U$ acting on the system and the bath copy and with $B=VI-[\Phi^*(H)-H+\Gamma^*(H_B)]$, a square unitary on systems the device holds during the use; the lift and pump above are the module's own whole-history phase lift and sharp pumps, which make it conserve energy and complete it on a finite box, so the module's bounds also return the phase reference; and the contact, which is the use's terminal contact, makes the energy gained by the system, the bath copy and the heat qubit equal to $w_c$ plus a remainder of norm at most $\|R_M\|$, uniformly over observers, a remainder that vanishes as $M\to\infty$. In the bits of this part, and since $H_B=-\beta^{-1}\ln\zeta_\delta$, it gives the nominal work per use
\[
\frac{\ln2}{\beta}\,D(\tau\|\zeta_\delta)+w_c+o(1)=\frac{\ln2}{\beta}\,g-\frac1\beta\ln(1-\delta)+o(1).
\]
It also gives everything else Theorem~\ref{sw:thm:clock} asks for: complete error $o(1)$ with respect to $n$ uses of the payload channel $\Phi_M$ of the primitive, which is within $\tfrac32b/M$ of $\Phi$ in half diamond norm, so the error with respect to $\Phi$ is $o(1)$ once $M\ge n^2b$ (unlike in Appendix~\ref{app:module}, the rotation here acts on every use rather than on a rare branch); a joint state of the whole controller with all other physical systems within $o(1)$ of the product; actual changes $o(1)$ of the controller's mean energy and entropy; and a battery that starts in an energy eigenstate and ends in a fixed energy eigenstate except with vanishing probability, with entropy and mean absolute deviation from that level $o(1)$, so that the work is deterministic. Retired bath copies and heat qubits stay in the reservoir. Lemma~\ref{sw:lem:minimax} and~\eqref{sw:eq:pointwise} hold for every channel, so $g$ can be chosen within any $a>0$ of $\Trt(\Phi)$. Letting $a,\delta\to0$ slowly with $n$ gives the rate $\kB T\ln2\,\Trt(\Phi)$.

\subsubsection{The range of the baths}

Suppose an implementation of $n$ uses of a unitary $U$ has error $\varepsilon$. For orthogonal inputs $a,b$, let $\sigma_a,\sigma_b$ be the final states of everything except the outputs. Feeding the superposition $(|0\rangle a+|1\rangle b)/\sqrt2$ with a kept qubit, and using that the ideal outputs are pure, the device's state must nearly factor off, which gives $T(\sigma_a,\sigma_b)\le2\sqrt{2\varepsilon}$~\cite{KSW08}. The argument follows Chiribella, Yang and Renner, who bound the capacity of a battery that implements a gate in the same way~\cite{ChiribellaYangRenner21}. Take $a$ and $b$ to be $n$-fold products of eigenvectors of $D$ for its largest and smallest eigenvalues. The delivered energies differ by at least $n(\lambda_{\max}-\lambda_{\min})(D)-2n\varepsilon\,\mathrm{span}(H)$. The work is deterministic, so the battery's energy changes by the same amount in both runs up to $o(n)$, and the controller's mean energy changes by $o(n)$ in both runs, so by energy conservation the mean energies of the heat baths differ by the same amount up to $o(n)$. For any bath Hamiltonian of total span $R$, $|\Tr H_{\mathrm{baths}}(\sigma_a-\sigma_b)|\le R\,T(\sigma_a,\sigma_b)$. Hence
\[
R\ \ge\ \frac{n(\lambda_{\max}-\lambda_{\min})(D)-2n\varepsilon\,\mathrm{span}(H)-o(n)}{2\sqrt{2\varepsilon}},
\]
which grows faster than $n$ as $\varepsilon\to0$ whenever $U$ is not time-covariant. Rare, large heat transfers are therefore unavoidable, as in the construction.

\subsection{Invertible environments}\label{sw:app:invertible}

We use the continuity bound of Alicki, Fannes and Winter~\cite{Winter16}: if two states of $XY$ satisfy $\tfrac12\|\rho-\sigma\|_1\le t\le1$, then
\begin{equation}\label{sw:eq:afw}
|S(X|Y)_\rho-S(X|Y)_\sigma|\le g_{\dim X}(t),\qquad g_D(t)=2t\log_2 D+(1+t)\,h\bigl(t/(1+t)\bigr),
\end{equation}
whatever the dimension of $Y$. We also use the stabilized norm $\|\Lambda\|_\diamond=\sup_k\|\Lambda\otimes\mathrm{id}_k\|_{1\to1}$, which is finite for every linear map $\Lambda$ between finite-dimensional matrix algebras.

\begin{proof}[Proof of Lemma~\ref{sw:lem:invertible}]
Let $\psi_{QOE}=(\mathrm{id}_Q\otimes V)\psi_\rho$, so that $\chi_\rho=\psi_{QO}$ and $\psi_{QE}=(\mathrm{id}_Q\otimes\Phi^c)(\psi_\rho)$. The state $\chi_\rho\otimes\sigma_F$ is feasible for every $\varepsilon$ and has $S(QO|F)=S(\chi_\rho)=S(\Phi^c(\rho))$, so $s_{\rho,\varepsilon}(\Phi)\le S(\Phi^c(\rho))$.

For the reverse inequality, let $\sigma_{QOF}$ be feasible at error $\varepsilon$ and purify it to $|\sigma\rangle_{QOFG}$. By the Fuchs--van de Graaf inequality the fidelity of $\sigma_{QO}$ and $\chi_\rho$ is at least $1-\varepsilon$. By Uhlmann's theorem there is an isometry $W:E\to FG$, enlarging $G$ if necessary, such that $|\langle\sigma|(I\otimes W)|\psi\rangle|\ge1-\varepsilon$; the two pure states are then within trace distance $\sqrt{2\varepsilon}$. Let $\Lambda(X)=\Tr_G WXW^*$ and $\tau_{QOF}=(\mathrm{id}_{QO}\otimes\Lambda)(\psi_{QOE})$. Then $\tfrac12\|\tau_{QOF}-\sigma_{QOF}\|_1\le\sqrt{2\varepsilon}$, and in particular $\tfrac12\|\tau_{QF}-\rho^T\otimes\sigma_F\|_1\le\sqrt{2\varepsilon}$.

Let $\mathcal R(X)=\Tr(X)\,\sigma_F$ and $\Delta=(\Lambda-\mathcal R)\circ\Phi^c$. Then $\tau_{QF}-\rho^T\otimes\sigma_F=(\mathrm{id}_Q\otimes\Delta)(\psi_\rho)$, whose trace norm is at most $2\sqrt{2\varepsilon}$. Since $\Delta$ preserves Hermiticity, its stabilized norm is attained on pure inputs with a reference of the same dimension~\cite[Thm.~3.51]{Watrous18}. Every unit vector of $Q\otimes A$ has the form $(X\otimes I)|\psi_\rho\rangle$ with $\|X\|^2\le1/\lambda_{\min}(\rho)$, so $\|\Delta\|_\diamond\le2\sqrt{2\varepsilon}/\lambda_{\min}(\rho)$. Since $\Phi^c$ is invertible, $\Lambda-\mathcal R=\Delta\circ(\Phi^c)^{-1}$, and
\[
\|\tau_{QOF}-\chi_\rho\otimes\sigma_F\|_1=\|(\mathrm{id}_{QO}\otimes(\Lambda-\mathcal R))(\psi_{QOE})\|_1\le\frac{2\sqrt{2\varepsilon}\,M}{\lambda_{\min}(\rho)},\qquad M=\|(\Phi^c)^{-1}\|_\diamond.
\]
Hence $\tfrac12\|\sigma_{QOF}-\chi_\rho\otimes\sigma_F\|_1\le t(\varepsilon):=\sqrt{2\varepsilon}\,(1+M/\lambda_{\min}(\rho))$, and~\eqref{sw:eq:afw} on the system $QO$ gives
\[
S(QO|F)_\sigma\ \ge\ S(\Phi^c(\rho))-g_{d_Qd_O}\bigl(\min\{1,t(\varepsilon)\}\bigr).
\]
The correction does not depend on $F$ and tends to zero with $\varepsilon$, which proves $s_\rho(\Phi)=S(\Phi^c(\rho))$.
\end{proof}

\subsection{Classical channels}\label{sw:app:classical}

Let $\Phi(\rho)=\sum_{x,y}P(y|x)\langle x|\rho|x\rangle\,|y\rangle\langle y|$ in an eigenbasis of $H$. Throughout, a \emph{dilation} is a finite unitary $U$ on the input and a bath in state $\tau$, possibly singular, with $\Phi(\rho)=\Tr_BU(\rho\otimes\tau)U^*$. Its bath output on input $\rho$ is $B'$, and $f_U(\rho)=S(B')-S(\tau)$. A dilation gives an exact extension in Definition~\ref{sw:def:rtcost}, with $F$ a purifier of $\tau$, and its value is $f_U(\rho)$; so $s_\rho(\Phi)\le f_U(\rho)$ for every dilation. Conversely $s_\rho(\Phi)\ge S(\rho)-S(\Phi(\rho))$ (Section~\ref{sw:sec:setting}).

\begin{lemma}[Coherence is free]\label{sw:lem:dephase}
Let $U$, with bath state $\tau$, implement $\Phi$ on diagonal inputs. Then there is a dilation $U'$ of $\Phi$ with $f_{U'}(\rho)\le f_U(\mathrm{diag}\,\rho)$ for every $\rho$.
\end{lemma}

\begin{proof}
Add to the bath a register $J$ in the maximally mixed state on $\mathbb Z_d$. First apply $\sum_j Z^j\otimes|j\rangle\langle j|_J$, where $Z=\sum_x e^{2\pi ix/d}|x\rangle\langle x|$, then $U$, which does not act on $J$. The input marginal is now $\mathrm{diag}\,\rho$, and by linearity the output and the bath output $B'$ are those of $U$ on $\mathrm{diag}\,\rho$; in particular $U'$ implements $\Phi(\mathrm{diag}\,\rho)=\Phi(\rho)$. By subadditivity, $S(J'B')\le\log d+S(B')$. Subtracting $S(J)+S(\tau)=\log d+S(\tau)$ gives the claim.
\end{proof}

\begin{lemma}[A nearly flat bath]\label{sw:lem:loguniform}
Let $\mu_N(l)=1/(l\,H_N)$ on $\{1,\dots,N\}$, with $H_N=\sum_{l\le N}1/l$. For every probability vector $w$ on a set of size $k$, list the products $\mu_N(l)\,w_s$ in decreasing order as $a_1\ge a_2\ge\cdots$. Then $a_r\le\mu_N(r)$ for $r\le N$, and
\[
\sum_{r>N}a_r\ \le\ \frac{1+\ln k}{\ln N}.
\]
\end{lemma}

\begin{proof}
The number of products exceeding $1/(rH_N)$ is $\sum_s\#\{l\ge1:l<rw_s\}<\sum_s rw_s=r$, so $a_r\le1/(rH_N)$. The $N$ largest products carry at least the mass of the products with $l\le\lfloor Nw_s\rfloor$, of which there are at most $N$. That mass is $\sum_sw_sH_{\lfloor Nw_s\rfloor}/H_N\ge\sum_sw_s\ln(Nw_s)/H_N=(\ln N-H_e(w))/H_N$, where $H_e$ is the entropy in nats; the inequality also holds for terms with $Nw_s<1$. Since $H_N\le1+\ln N$ and $H_e(w)\le\ln k$, the remaining mass is at most $(1+\ln k)/H_N\le(1+\ln k)/\ln N$.
\end{proof}

\begin{proof}[Proof of Theorem~\ref{sw:thm:classical}(a)]
\emph{Diagonal inputs.} Fix a full-rank probability vector $p$. The rows $P(\cdot|x)$ have at most $d(d-1)$ interior breakpoints between them in their cumulative distributions, so there is a random variable $R$ with at most $d^2$ values, independent of $X$, and a function $f$ with $Y=f(X,R)$ distributed exactly as $P(\cdot|X)$; the probabilities of $R$ are interval lengths and may be irrational. For each $y$ let $\mathcal F_y=f^{-1}(y)$ and let $w_y$ be the conditional law of $(X,R)$ given $Y=y$. The bath holds $R$, a register $L\sim\mu_N$ and a padding register in a fixed state. The dilation is a permutation of classical labels that sends $(x,r,l)$ to $(y,b)$ with $y=f(x,r)$ and $b$ the rank of $\mu_N(l)\,w_y(x,r)$ among the products for that $y$, completed arbitrarily on labels of weight zero. It implements $\Phi$ exactly on diagonal inputs, and Lemma~\ref{sw:lem:dephase} turns it into a dilation. Each fibre has $|\mathcal F_y|\le d^3$ elements, so by Lemma~\ref{sw:lem:loguniform} the conditional law of $B'$ given $Y=y$ is within total variation $(1+3\ln d)/\ln N$ of $\mu_N$, extended by zeros, for every $y$. Comparing the joint law of $YB'$ with the product of the law of $Y$ and $\mu_N$, and using~\eqref{sw:eq:afw} on $Y$, gives $I(Y{:}B')\le g_d\bigl((1+3\ln d)/\ln N\bigr)$. Since a permutation conserves entropy, $f_U(p)=H(p)-H(Pp)+I(Y{:}B')$. Letting $N\to\infty$ gives $s_p(\Phi)\le H(p)-H(Pp)$, hence equality.

\emph{All inputs.} By Lemma~\ref{sw:lem:dephase}, $s_\rho(\Phi)\le H(\mathrm{diag}\,\rho)-H(P\,\mathrm{diag}\,\rho)$. The energy term $\Tr H(\Phi(\rho)-\rho)$ depends only on $\mathrm{diag}\,\rho$, because $H$ is diagonal and $\Phi(\rho)=\Phi(\mathrm{diag}\,\rho)$. So the supremum in~\eqref{sw:eq:realtime} is at most the maximum in~\eqref{sw:eq:batch} over diagonal inputs, which is at most $\Tcap(\Phi)$. Since $\Trt\ge\Tcap$ always, they are equal. A non-full-rank diagonal input is a limit of full-rank ones, and the batch maximand is continuous.
\end{proof}

The second step is needed: on a coherent input the correlation term of a classical channel can be positive, so the proof bounds the optimized cost rather than $m_\rho$. Take complete dephasing $\Delta$ of a qubit in the computational basis, with $H=0$. Its basis outputs are pure, so every dilation acts on the support of $\tau$ as $|i\rangle|v\rangle\mapsto|i\rangle V_i|v\rangle$ with isometries $V_i$, and on input $\rho$ the bath ends in $\sum_i\rho_{ii}V_i\tau V_i^*$, a mixture of states with the spectrum of $\tau$. By concavity its entropy is at least $S(\tau)$, so $s_\rho(\Delta)\ge0$, with equality for a controlled $Z$ and a maximally mixed qubit bath. Hence $m_\rho(\Delta)=S(\Delta(\rho))-S(\rho)$, which for $\rho=(I+0.6X)/2$ is $1-h(0.8)\approx0.278$ bits, for every dilation and not only for a poor one. At that input the batch integrand is $-0.278$ and the real-time integrand is $0$; both maxima are $0$, attained on diagonal inputs, where $m_p=0$.

The device depends on the input $p$, through the ranking. This is unavoidable: for the binary channel with $P(1|1)=1/2$ and $P(0|0)=1$, no single dilation leaves its bath nearly independent of the output for every input. Theorem~\ref{sw:thm:law} does not need such a device, because $\Trt$ is defined input by input and the construction of Appendix~\ref{sw:app:construction} combines input-dependent dilations through a minimax argument.

\begin{proof}[Proof of the attainment remark]
Here the Hamiltonian is zero, so $\Trt=\Tcap=\max_p[H(p)-H(Pp)]$. For unital $P$, Birkhoff's theorem writes $P$ as a mixture of permutations $\Pi_k$, so $\Phi$ is the mixture of the unitaries $\Pi_kZ^j$ with weights $\lambda_k/d$. The controlled unitary with a diagonal bath carrying these weights leaves the bath marginal unchanged, so $f_U=0=\Tcap$. For deterministic $P$, given by a function $y=f(x)$, the map $(x,0)\mapsto(f(x),\text{the index of }x\text{ in }f^{-1}(f(x)))$, preceded by the register $J$ of Lemma~\ref{sw:lem:dephase}, has $f_U\le\log_2\max_y|f^{-1}(y)|=\Tcap$.

Now let $P(0|0)=1$ and $P(0|1)=a\in(0,1)$, and suppose a finite dilation $U$ with bath spectrum $t$ attains $\Tcap$. Its cost on a diagonal input $p$ is $H(p)-H(Pp)+I(Y{:}B')$, so attainment forces $I(Y{:}B')=0$ at the batch-optimal input $p^*$, which has full rank. Let $w_x=U(|x\rangle\langle x|\otimes\tau)U^*$; both have spectrum $t$, and they are orthogonal. Since input $0$ is sent to $0$, $w_0=|0\rangle\langle0|\otimes\theta_0$. Independence at $p^*$ means $p_0^*w_0+p_1^*w_1=\mathrm{diag}(q_0,q_1)\otimes\omega$ with $q=Pp^*$. The $|1\rangle$ block gives $\langle1|w_1|1\rangle=(1-a)\,\omega$, and the off-diagonal block gives that $w_1$ is block diagonal, so $t\supseteq(1-a)\,\mathrm{spec}(\omega)$. The $|0\rangle$ block gives $q_0\,\omega=p_0^*\theta_0+p_1^*\langle0|w_1|0\rangle$ with orthogonal supports, so $\mathrm{spec}(\omega)\supseteq(p_0^*/q_0)\,t$. The smallest positive element $m$ of $t$ then satisfies $m\le(1-a)(p_0^*/q_0)\,m<m$, a contradiction.
\end{proof}

\subsection{The Werner--Holevo support test}\label{sw:app:werner}

Let $P_-$ be the projector onto the antisymmetric subspace of $QO$, with $Q$ and $O$ qutrits, and $V$ the isometry of Section~\ref{sw:sec:werner}, so that $P_-=VV^*$ and $VI_3V^*/3=P_-/3$ is the Choi state $J_W$ of $W$.

\begin{proposition}\label{sw:prop:support}
In a real-time implementation of $W$ with Bell-pair inputs, let $\delta_t$ be the probability that $Q_tO_t$ fails the test $\{P_-,I-P_-\}$. Then the entropy of the device rises in round $t$ by at least $\bigl[\log_2 3-g_9\bigl(\min\{1,10\sqrt{\delta_t}\}\bigr)\bigr]_+$, and the free energy drawn satisfies $w\ge n\bigl[\log_2 3-g_9(\min\{1,10\sqrt{\bar\delta}\})\bigr]_+$ with $\bar\delta=\frac1n\sum_t\delta_t$.
\end{proposition}

\begin{proof}
Let $\sigma=\sigma_{QOF}$ be the state after round $t$, with $F$ as in Lemma~\ref{sw:lem:accumulation}, so that $\sigma_{QF}=\tau_3\otimes\beta_F$ with $\tau_3=I_3/3$. Project onto $P_-\otimes I_F$ and renormalize to get $\sigma'$; the gentle-measurement lemma gives $\|\sigma-\sigma'\|_1\le2\sqrt{\delta_t}$. Since $\sigma'$ is supported on $\mathrm{ran}(V)\otimes F$, it equals $(V\otimes I)\zeta(V^*\otimes I)$ for a state $\zeta$ on $\mathbb C^3\otimes F$, and $\sigma'_{QF}=(W\otimes\mathrm{id})(\zeta)$ because tracing out $O$ after $V$ is $W$. The inverse $W^{-1}(Y)=\Tr(Y)I-2Y^T$ has $\|W^{-1}\|_\diamond\le3+2\cdot3=9$, and $W^{-1}(\tau_3)=\tau_3$. Applying it,
\[
\|\zeta-\tau_3\otimes\beta_F\|_1\le9\,\|\sigma'_{QF}-\tau_3\otimes\beta_F\|_1\le9\cdot2\sqrt{\delta_t}.
\]
Conjugating by $V$, which preserves the trace norm and maps $\tau_3\otimes\beta_F$ to $J_W\otimes\beta_F$, gives $\|\sigma-J_W\otimes\beta_F\|_1\le20\sqrt{\delta_t}$. By~\eqref{sw:eq:afw} on the nine-dimensional system $QO$,
\[
S(QO|F)_\sigma\ge S(J_W)-g_9(\min\{1,10\sqrt{\delta_t}\})=\log_23-g_9(\min\{1,10\sqrt{\delta_t}\}).
\] The increment is also nonnegative, since $S(QO|F)\ge S(Q)-S(O)\ge0$ for a maximally mixed $Q$. Summing over rounds and using Lemma~\ref{sw:lem:accounting} gives the first form. The function $\delta\mapsto g_9(\min\{1,10\sqrt\delta\})$ is increasing and concave on $[0,1]$, so $\delta\mapsto[\log_23-g_9(\min\{1,10\sqrt\delta\})]_+$ is convex and Jensen's inequality gives the averaged form.
\end{proof}

\subsection{Zero error and extreme channels}\label{sw:app:extreme}

\begin{proof}[Proof of Proposition~\ref{sw:prop:extreme}]
Let the device start in $\tau_D=\sum_a\lambda_a|a\rangle\langle a|$, independent of the inputs, and let it implement $\Phi^{\otimes n}$ exactly. Each $|a\rangle$ gives an isometric dilation $U_a$ of a channel $\Phi_a$ with $\sum_a\lambda_a\Phi_a=\Phi^{\otimes n}$; since $\Phi^{\otimes n}$ is extreme, $\Phi_a=\Phi^{\otimes n}$ for every $a$. By uniqueness of the minimal dilation $V$, $U_a=(I\otimes T_a)V$ for isometries $T_a$, and unitarity gives $U_a^*U_b=\delta_{ab}I$, that is $V^*(I\otimes T_a^*T_b)V=\delta_{ab}I$. Extremality makes $X\mapsto V^*(I\otimes X)V$ injective, so $T_a^*T_b=\delta_{ab}I$. The final device state is therefore the orthogonal sum $\sum_a\lambda_aT_a\Phi^{c\otimes n}(\rho)T_a^*$, and its entropy rises by exactly $S(\Phi^{c\otimes n}(\rho))$ on every input. Lemma~\ref{sw:lem:accounting} then gives $w\ge A+\beta\Delta E_{\mathrm{sys}}/\ln2=n\bigl[S(\Phi^c(\rho))+\tfrac{\beta}{\ln2}\Tr H(\Phi(\rho)-\rho)\bigr]$ on $\rho^{\otimes n}$, with no error term; the $o(n)$ enters only when $w$ is converted into mechanical work, using that the battery's entropy changes by $o(n)$ and that the controller returns. Entropy is subadditive, so entangled inputs give no more, and the bound is $n\,a_\beta(\Phi)$.

It remains to show $s_\rho(\Phi)=S(\Phi^c(\rho))$ for full-rank $\rho$; the minimal dilation gives $\le$. Let $\sigma_{QOF}$ be feasible in Definition~\ref{sw:def:rtcost} at error $\varepsilon$. A rank-one measurement on $F$ leaves $Q$ in $\rho^T$, so it splits $\sigma_{QO}$ into states $(\mathrm{id}\otimes\Phi_k)(\psi_\rho)$ of channels $\Phi_k$ whose average is within $O_\rho(\varepsilon)$ of $\Phi$. Because $\Phi$ is extreme, the point mass at $\Phi$ is the only probability measure on channels with barycentre $\Phi$, and compactness gives $\sum_kq_k\|\Phi_k-\Phi\|_\diamond\le g_\rho(\varepsilon)\to0$ uniformly over all such splittings. Expanding $\sigma_{QOF}-\sigma_{QO}\otimes\sigma_F$ in an orthonormal operator basis of $QO$ bounds its trace norm by $2d^5(g_\rho(\varepsilon)+\varepsilon)$, whatever the size of $F$, and the Alicki--Fannes--Winter inequality~\cite{Winter16}, whose error depends only on $\dim QO=d^2$, gives $S(QO|F)_\sigma\ge S(\chi_\rho)-f_d(\varepsilon)$ with $f_d(\varepsilon)\to0$. Since $S(\chi_\rho)=S(\Phi^c(\rho))$, we get $s_\rho(\Phi)=S(\Phi^c(\rho))$, and hence $a_\beta(\Phi)=\Trt(\Phi)$. The modulus $f_d$ is steep, so finite-size statements at small error must carry it.
\end{proof}

% =====================================================================
\subsection{Exact inputs and the fixed-dimension reduction}\label{sw:app:fixeddim}
This appendix proves Theorem~\ref{sw:thm:undecidable} in full. We fix exact encodings of the inputs, state the results we use, prove a transfer lemma from correlations to Gram matrices of unitaries, and assemble the reduction. Every test is exact; no step uses numerical approximation. Let
\[
\mathbb K_0=\bigcup_{\ell\ge1}\mathbb Q(\zeta_\ell),\qquad \mathbb K=\mathbb K_0\cap\mathbb R.
\]
Elements have finite exact encodings, for example rational polynomials modulo a cyclotomic polynomial, with $\zeta_\ell=e^{2\pi i/\ell}$. Bring any finite collection to one common field using the least common multiple of root orders. Equality, conjugation, field arithmetic and signs of real algebraic expressions are decidable; isolating-real-root algorithms give the signs. Gaussian elimination gives exact ranks, and a Hermitian matrix is PSD exactly when all its principal minors, which are real, are nonnegative. These tests are necessary conditions only; they do not decide membership in a closure.

\paragraph{Notation.} For $N$ questions and eight answers, $C_{qc}^s$ and $C_{qa}^s$ are the synchronous correlations of commuting-operator and approximately quantum type. $\G_{\rm matr}(m)$ is the set of $m\times m$ Gram matrices $[\tau(u_j^*u_i)]$ of unitaries in full matrix algebras with normalized trace. $\mathrm{FM}_{\rm matr}(d)$ and $\mathrm{FM}_{\rm fin}(d)$ are the unital Schur channels on $M_d$ that factor through full matrix algebras and through finite-dimensional algebras with a faithful trace; $\K_d$ is the closure of either set. $\R$ is the hyperfinite II$_1$ factor and $\R^\omega$ its ultrapower, not the reservoir of Section~\ref{sw:sec:setting}.

\paragraph{Results we use.}
Fu, Miller and Slofstra~\cite{FMS}, in the opening paragraph of their Section~7, replace the earlier field of all real algebraic numbers by $\mathbb K=\mathbb K_0\cap\mathbb R$. Their Theorem~7.2 gives, for any recursively enumerable set $X$, a fixed $N$ and computable finite lists $F_n\subset\mathbb K^{N^2\times8^2}$ such that $n\in X$ gives no $C_{qc}$ member and $n\notin X$ gives a $C_{qa}$ member; the construction is synchronous. Section 6, equation (2), uses rational coefficients and sines/cosines of integer multiples of $\pi/p$. Section 7's $\sigma(x,a)$ are finite products of those projections and rational group-algebra expressions; $C_f=f(\sigma(x,a)\sigma(y,b))$ uses finite coefficient sums with $f$ valued in $\{0,1\}$. Thus each list is computably cyclotomic, not merely algebraic. Indeed sine/cosine coefficients lie in $\mathbb Q(\zeta_{4p})$, and adding the transfer's $i$ and $\sqrt2$ stays in $\mathbb Q(\zeta_{8p})$. Root orders and degrees grow with $n$; no one fixed number field is claimed. The choice of a suitable primitive root and the resulting $N$ is existential, so the theorem does not yield a numerical dimension here.

Kim, Paulsen and Schafhauser~\cite[Thm.~3.6]{KPS18} identify synchronous $C_{qa}$ with limits of finite-dimensional synchronous tracial projection correlations and with projection correlations in $\R^\omega$. Musat and R{\o}rdam~\cite[Prop.~3.2]{MusatRordam} supply the projection-to-unitary trick. Musat~\cite[Prop.~3.9(iii)]{Musat25} identifies a Schur channel in the closure of matrix-factorizable channels exactly with its matrix in $\overline{\G_{\rm matr}}$. For completeness we prove the transfer.

\begin{lemma}[Formal faithful Gram transfer]\label{sw:fd-lem:gram}
Let $P(a,b|x,y)$ be a real synchronous normalized table with $N$ questions and eight answers, with consistent marginals. Enumerate $m=8N$ symbols $p_j=p_{x,a}$. Put
\[
M_{jk}=P(a,b|x,y),\quad q_j=P(a,a|x,x),\quad
A=\begin{pmatrix}1&q^T\\q&M\end{pmatrix}.
\]
Define the formal vectors
\[
u_0=1,\qquad u_j=2p_j-1,\qquad
u_{m+j}=(u_j+i1)/\sqrt2\quad(1\le j\le m),
\]
and let $C_P$ be their formal Gram $[\tau(u_j^*u_i)]$, computed from $A$ ($i,j=0,\dots,2m$; by traciality this equals $[\tau(u_iu_j^*)]$, the convention of~\cite{Musat25}). For tables that pass the exact consistency and PSD filters,
\begin{equation}\label{sw:fd-eq:equivalence}
P\in C_{qa}^{s}\quad\Longleftrightarrow\quad
C_P\in\overline{\G_{\rm matr}(2m+1)}.
\end{equation}
If $C_P$ is a unitary Gram in any faithful finite tracial von Neumann algebra, then $P\in C_{qc}^{s}$.
\end{lemma}
\begin{proof}
All entries of $C_P$ are affine expressions in $P$ with coefficients in $\mathbb Q(i,\sqrt2)$; it is computed by the coefficient matrix acting on $A$. A tracial PVM representation of $P$ makes the displayed $u_j$ actual unitaries and gives this Gram. By the theorem of Kim, Paulsen and Schafhauser, if $P\in C_{qa}^s$, finite-dimensional synchronous PVM correlations approach it. The same formulas give approaching unitary Grams. If their traces use finite direct sums, approximate the weights by rational weights and duplicate blocks, giving full matrix normalized traces with arbitrarily small moment error. This proves the forward implication.

For the converse first work in an arbitrary faithful finite tracial algebra realizing $C_P$ by unitaries $v_i$. Replace all of them by $v_0^*v_i$ so $v_0=1$, without changing any Gram entry. The formal relation
\[
u_{m+j}-(u_j+i u_0)/\sqrt2=0
\]
is a null vector of the Gram. Its corresponding actual $L^2(\tau)$ norm is zero, and faithfulness implies
$v_{m+j}=(v_j+i1)/\sqrt2$. Since both are unitary,
\[
1=v_{m+j}^*v_{m+j}=I+\tfrac i2(v_j^*-v_j),
\]
so $v_j=v_j^*$ and $r_j=(v_j+I)/2$ is a projection. Expanding the known Gram gives $\tau(r_j)=q_j$ and $\tau(r_kr_j)=M_{jk}$. For each fixed question, synchronicity gives $\tau(r_{x,b}r_{x,a})=0$ for $a\ne b$. But $\|r_{x,b}r_{x,a}\|_2^2=\tau(r_{x,a}r_{x,b})=0$, so these projections are orthogonal. Their sum is a projection of trace $\sum_aq_{x,a}=1$, hence equals $I$ by faithfulness. They are PVMs with correlations $P$.

Such a tracial PVM representation yields a commuting correlation: in $L^2(\mathcal A,\tau)$, left multiplication by $r_{x,a}$ and right multiplication by $r_{y,b}$ commute and, on vector $1$, have correlation $\tau(r_{x,a}r_{y,b})$. This proves the last assertion.

If $C_P$ is a limit of full matrix unitary Grams, take the tracial ultraproduct of their matrix algebras. The coordinatewise unitaries have exact Gram $C_P$ in the faithful ultraproduct trace. Each matrix embeds trace-preservingly into the hyperfinite factor $\R$; these embeddings induce an embedding of the ultraproduct into $\R^\omega$. The preceding null-vector argument gives exact PVMs in $\R^\omega$ realizing $P$. That theorem therefore gives $P\in C_{qa}^s$. This proves the reverse implication.
\end{proof}

\begin{theorem}[Fixed dimension, finite-list non-enumerability]\label{sw:fd-thm:fixed}
Fix an RE-complete set $X$, and let $N$ be as in~\cite[Thm.~7.2]{FMS}. Set $d_0=16N+1$. For every $d\ge d_0$ there is a total computable map $n\mapsto L_n$ to a finite list of unital Schur channels on $M_d$ with exact entries in $\mathbb K_0$ such that
\begin{align*}
n\notin X&\quad\Longleftrightarrow\quad
\text{some }\Sigma\in L_n\text{ lies in }\K_d,\\
n\in X&\quad\Longrightarrow\quad
\text{no }\Sigma\in L_n\text{ is factorizable},
\end{align*}
where $\K_d=\overline{\mathrm{FM}_{\rm matr}(d)}=\overline{\mathrm{FM}_{\rm fin}(d)}$. Membership in $\K_d$ on these exact Schur inputs is not recursively enumerable, and hence undecidable. The reduction is a disjunctive truth-table reduction, and $d_0$ exists but is not given numerically.
\end{theorem}
\begin{proof}
Compute $F_n$. For each table make exact tests of realness, nonnegativity, normalization, synchronicity, marginal consistency, the required symmetry of $M$, PSD of $A$, and PSD/unit diagonal of $C_P$. Discard any failure. Every $C_{qa}$ table passes these necessary conditions, so no positive-side witness is lost; all the tests use exact arithmetic.

For each surviving table output $\Sigma_{C_P}(E_{ij})=(C_P)_{ij}E_{ij}$. PSD and unit diagonal prove complete positivity, trace preservation and unitality. The size is $2(8N)+1=d_0$. The entries stay in $\mathbb K_0$ because the formal transfer only adds $i$ and $\sqrt2$. If $n\notin X$, some table of $F_n$ is synchronous $C_{qa}$; Lemma~\ref{sw:fd-lem:gram} and~\cite[Prop.~3.9(iii)]{Musat25} put its channel in $\K_{d_0}$. If $n\in X$, none of the input tables is $C_{qc}$. Any factorizable output Schur channel would, by the unitary-Gram characterization \cite[Prop.~3.9(i)]{Musat25}, give a Gram in a faithful finite tracial algebra; the final assertion of Lemma~\ref{sw:fd-lem:gram} would put its table in $C_{qc}$, contradiction. This is stronger than merely lying outside $\K_{d_0}$.

For $d>d_0$ choose a surjection $f:\{1,\ldots,d\}\to\{1,\ldots,d_0\}$ equal to the identity on the first $d_0$ indices, and put $C'_{ij}=C_{f(i),f(j)}$. Duplicating Gram vectors proves PSD and unit diagonal. A unitary representation or an approximating sequence extends by duplicating the corresponding unitaries; conversely, restricting to the first $d_0$ indices recovers the old Gram. The closure and factorizability predicates are therefore both preserved in every larger dimension. All padded entries remain exact cyclotomic.

If $\K_d$ membership were RE on these exact inputs, enumerate accepted channels and dovetail the enumerator on the finite computed list $L_n$. Accept $n$ as soon as any member is accepted. This would enumerate $\mathbb N\setminus X$, contradicting the RE-completeness and nonrecursiveness of $X$. An empty list simply never accepts. If a decider existed, evaluate all finitely many members and OR the answers, deciding $\mathbb N\setminus X$. This is a finite disjunction of nonadaptive membership queries.
\end{proof}

\paragraph{From factorizability to work.}
For a Schur channel and diagonal $H$, every input population and its mean energy are preserved, so the energy term of the real-time functional is zero for every input. So $\Trt$ is the zero-Hamiltonian purity cost, and by the zero-Hamiltonian closure theorem of~\cite[Thm.~3.2(a)]{Memory} and~\cite{Douglas26closed} it vanishes exactly when the channel belongs to $\K_d$. As in Step~1 of Section~\ref{sw:sec:undecidable}, the hard direction of Theorem~\ref{sw:thm:free} is not used, and nothing here comes from~\cite{FMS}. Its easy thermal realization can be checked directly: if
$\mathcal A=\bigoplus_kM_{n_k}$ and $\tau=\sum_kw_k\tr_{n_k}$, choose
\[
e^{-\beta E_k}=w_k/n_k,\qquad H_B=\bigoplus_kE_kI_{n_k},\qquad
U=\sum_i|i\rangle\langle i|\otimes u_i.
\]
Then the bath partition function is one, its block density is $w_kI/n_k$, and $U$ commutes with diagonal payload $H$ and the block-scalar $H_B$. It gives the prescribed Schur multiplier exactly; blocks of zero weight are removed, and with unnormalized traces the factor $n_k$ must be kept. Limits stay free, and zero cost forces membership in $\K_d$, by the zero-Hamiltonian closure theorem~\cite[Thm.~3.2(a)]{Memory} (the first direction also~\cite[Sec.~8]{Douglas26closed}). Together with Theorem~\ref{sw:fd-thm:fixed}, this proves Theorem~\ref{sw:thm:undecidable} for every fixed diagonal $H$ and every $\beta>0$.

\paragraph{Scope.}
The dimension $d_0$ is not computed, and the cyclotomic fields grow with $n$. The reduction is a disjunctive truth-table reduction, not a many-one reduction, and it gives no uniform distance from $\K_d$ on the hard side, where the channels are not factorizable. It covers neither rational inputs nor devices with a fixed delay. Exact membership in $\K_d$ is in $\Pi^0_2$: for every precision there exists a matrix size and unitary tuple achieving it, and each fixed-size semialgebraic feasibility predicate over exact algebraic coefficients is decidable. This does not imply RE or decidability.

% =====================================================================
